\section{Strategy, Backtesting, and Evaluation Framework}
\label{sec:strategy}

\subsection{Strategy Archive}

The repository preserves the full evolution of the strategy line whose failure
motivated this release (v1--v15 iterations), providing historical baselines
for comparison and reproducibility. Strategies combine model confidence,
market price, entry constraints, and risk controls. The current research line includes
drift and order-flow imbalance signals, scoreboard-derived Polymarket book
signals, whipsaw/chop detection, best-signal scanning inside an entry window,
volume gating, entry price bounds, confidence and edge thresholds, Brier-score
circuit breakers, and confidence-bin empirical edge gates.

Position sizing and simulated execution assumptions are intentionally separated
from directional accuracy. Backtests report slippage, fees, assumed fill
prices, and market impact assumptions explicitly.

\subsection{Backtesting}
\label{sec:backtesting}

Backtesting processes market windows independently and evaluates entry signals
against settled outcomes under documented simulation assumptions (signal gate,
simulated fill, settlement, metric aggregation). Preferred
validation methods include walk-forward validation, rolling-window evaluation,
strict out-of-sample splits, sensitivity analysis over entry windows and
thresholds, and Monte Carlo resampling of fill and slippage assumptions. Random
row splits are discouraged because adjacent high-frequency observations are
highly autocorrelated.

\subsection{Evaluation Metrics}
\label{sec:evaluation}

Evaluation includes both predictive and simulated-economic metrics.
Economic metrics are counterfactual diagnostics under explicit fill, fee, and
slippage assumptions.

\paragraph{Predictive metrics.}
Accuracy, precision, recall, F1, ROC AUC, Brier score, and calibration
curves~\cite{brier1950verification}.

\paragraph{Simulated-economic metrics.}
Simulated win rate, expectancy, PnL, Sharpe and Sortino ratios (under stated
assumptions), maximum drawdown, turnover, average entry price, and documented
fill, fee, and slippage assumptions.

Calibration is especially important because binary market strategies are
sensitive to the difference between predicted probability and market-implied
price.
